% macros-repaso.tex — Equivalencias en LaTeX llano de las órdenes de tcee-repaso.sty
% (y de las de tcee-comun.sty que usan las fichas) para que pandoc las convierta a
% Markdown. scripts/skill-opositor.py lo antepone al cuerpo de cada ficha de repaso
% (fuentes/<T>/repaso/<T>-repaso.tex) y del resumen de la guía del dictamen (1D01).
%
% Lo que una macro no puede calcular lo sustituye antes el script en Python: los
% minutos del cante (\bloque[2,5]{…} → \bloquemin{2′30″}{…}; el segundo argumento de
% \apartadocante y \subcante pasa a «2′30″», o a \apartadocantesin{…} y \subcantesin{…}
% si no hay minutos; \tiempo{2,5} → 2′30″) y el título de la caja «trampa»
% (\trampatitulo: «Ojo en el test» en las fichas, «Ojo en el dictamen» en el resumen).
%
% Si se añade una orden nueva a tcee-repaso.sty, hay que añadir aquí su equivalente.
% Autor: Víctor Gutiérrez Marcos

% --- Cabecera y estructura del documento (el script escribe el título en Markdown) ---
\newcommand{\tcetema}[3][]{}
\newcommand{\tcefecha}[1]{}
\newcommand{\cabecerarepaso}{}
\newcommand{\tceeresumen}{}
\newcommand{\tceetam}{}
\newcommand{\tceetitulo}{}
\newenvironment{columnas}{}{}
\providecommand{\trampatitulo}{Ojo en el test}

% --- Bloques y subapartados: títulos de nivel 1 y 2 (pandoc los baja un nivel) ---
\newcommand{\bloque}[2][]{\section{#2}}
\newcommand{\bloquemin}[2]{\section{#2 (#1)}}
\newcommand{\sub}[1]{\subsection{#1}}

% --- Esquema del cante: una lista con los minutos de cada apartado ---
\newcommand{\esquemacante}[1]{\begin{itemize}#1\end{itemize}}
\newcommand{\apartadocante}[2]{\item \textbf{#1} · #2}
\newcommand{\subcante}[2]{\item #1 · #2}
\newcommand{\apartadocantesin}[1]{\item \textbf{#1}}
\newcommand{\subcantesin}[1]{\item #1}
\newcommand{\hito}[2]{\par\textbf{#1} #2\par}

% --- Énfasis y remisiones ---
\newcommand{\autor}[1]{#1}
\newcommand{\concepto}[1]{\textbf{#1}}
\newcommand{\vertema}[1]{[ver tema #1]}
\newcommand{\vertemacorto}[1]{#1}
\newcommand{\dato}[2]{\textbf{#1} (#2)}
\newcommand{\minutos}[1]{(#1)}
\renewcommand{\hyperref}[2][]{#2}
\newcommand{\findemo}{∎}

% --- Gráficos: solo el pie, entre corchetes ---
\newcommand{\graficorepaso}[3][]{\par\emph{[Gráfico: #3]}\par}
\newcommand{\pareja}[4]{\par\emph{[Gráfico: #2]} · \emph{[Gráfico: #4]}\par}
\renewcommand{\includegraphics}[2][]{\emph{[Gráfico: #2]}}
\newcommand{\fuente}[1]{\par\emph{Fuente: #1}\par}

% --- Fórmulas ---
\newcommand{\funcionclave}[1]{\[#1\]}
\providecommand{\vect}[1]{\mathbf{#1}}
% \optim{\max}{x}{objetivo}{restricciones}: «max» con la variable debajo y «s.a.»
\newcommand{\optim}[4]{\begin{array}{rl}#1_{#2} & #3 \\ \text{s.a.} & #4\end{array}}

% --- Cajas: la fórmula se queda como está; idea clave y trampa, como citas ---
\newenvironment{formula}{}{}
\newenvironment{ideaclave}{\begin{quote}}{\end{quote}}
\newenvironment{trampa}{\begin{quote}\textbf{\trampatitulo.} }{\end{quote}}

% --- Tablas: lo que pandoc no entiende se vacía ---
\newcommand{\rowcolor}[1]{}
\newcommand{\cellcolor}[1]{}
\newcommand{\arraybackslash}{}
\newcommand{\addlinespace}[1][]{}
